% ~400 words
% Key: name the gaps, then give 3 explanations (documentation / mental model / workflow barrier)

Table~\ref{tab:gaps} summarises the most pronounced adoption gaps ---
features that were deliberately designed into LiaScript but are used
in fewer than \SI{5}{\percent} of courses.

\begin{table}[ht]
  \centering
  \caption{Underutilised LiaScript features (adoption $< \SI{5}{\percent}$)}
  \label{tab:gaps}
  \begin{tabular}{llrp{4.2cm}}
    \toprule
    \textbf{Feature} & \textbf{Category} & \textbf{Adoption} & \textbf{Design intent} \\
    \midrule
    Effects            & Presentation  & \SI{0.0}{\percent}  & Timed content appearance \\
    TTS blocks         & Presentation  & \SI{0.1}{\percent}  & Fine-grained audio control \\
    Code projects      & Embedding     & \SI{0.4}{\percent}  & Multi-file code environments \\
    Animated CSS       & Presentation  & \SI{0.7}{\percent}  & Visual step-by-step reveals \\
    Surveys            & Interaction   & \SI{2.7}{\percent}  & Learner feedback collection \\
    Exec.\ code        & Embedding     & \SI{3.1}{\percent}  & Executable code environments \\
    Task lists         & Interaction   & \SI{4.7}{\percent}  & Self-directed checklists \\
    \bottomrule
  \end{tabular}
\end{table}

We identify three classes of explanation for adoption gaps:

\paragraph{Documentation deficit.}
Some features are syntactically non-obvious and lack prominent examples.
TTS blocks, animated CSS, and effects fall into this category:
the basic narrator comment (\texttt{<!-- voice: ... -->}) is well documented,
but the richer block-level TTS and animation system requires authors to
understand a secondary syntax layer.
Better worked examples and interactive tutorials would likely raise adoption.

\paragraph{Genuine workflow barrier.}
Code projects (multi-file environments) and some embedding features require
authors to configure external services or understand JavaScript interop.
For educators without a technical background this represents a real barrier
that documentation alone cannot resolve.
Tool-level improvements --- wizards, visual editors, or pre-configured
project templates --- are needed.

The distribution across these two classes directly informs the
development priorities discussed in Section~\ref{sec:implications}.
